% !TEX program = xelatex
\documentclass[12pt]{article}
\usepackage[margin=28mm]{geometry}
\usepackage{hyperref}
\usepackage{amsmath, amsthm, amssymb, amsfonts}
\usepackage{tikz-cd}
\usepackage{fontspec}
\usepackage{unicode-math}
\usepackage{xeCJK} % 日本語
\setmainfont{Latin Modern Roman}
\setmonofont{Latin Modern Mono}
\setCJKmainfont{Noto Serif CJK JP} % お手元の環境に応じて変更可（例：IPAexMincho）
\usepackage{listings}
\lstset{basicstyle=\ttfamily\footnotesize, columns=fullflexible, breaklines=true, showstringspaces=false}
\newtheorem{theorem}{定理}
\newtheorem{lemma}{補題}
\newtheorem{proposition}{命題}
\theoremstyle{definition}
\newtheorem{definition}{定義}
\newtheorem{remark}{注意}

\title{BourbakiPairsAndMorphisms — 日本語数学解説付きノート}
\author{TeX generated by ChatGPT-5 Pro\\Lean source generated by CODEX (OpenAI)}
\date{2025-09-08}

\begin{document}
\maketitle

\begin{abstract}
本資料は Lean ファイル \texttt{BourbakiPairsAndMorphisms.lean} から自動生成し、
\textbf{著者（生成元）の明示}として
「TeXは ChatGPT-5 Pro により生成」「Leanファイルは CODEX (OpenAI) により生成」
であることを記載しています。Bourbaki 的な視点で、積位相、群準同型、位相群に関する
主要命題の\textbf{日本語による数学的解説}と\textbf{図式}を付しました。
\end{abstract}

\section*{Provenance（生成元の明示）}
\begin{itemize}
  \item \textbf{TeX generation}: ChatGPT-5 Pro
  \item \textbf{Lean source generation}: CODEX (OpenAI)
\end{itemize}

\section*{背景と狙い}
本ノートは、**Bourbaki 風の観点**（集合に担われた構造、射影 $\pi_1,\pi_2$ による順序対の特徴づけ、
および射の判定を射影に還元する考え方）に沿って、Lean ファイルから得られる命題を**日本語で解説**するものです。

題材は大きく次の 3 テーマに分かれます。

- **A. 積位相**：$h:T\to X\times Y$ の連続性を $\pi_1\circ h,\ \pi_2\circ h$ の連続性で特徴づける。
- **B. 群準同型の合成**：群（正確には MonoidHom）における合成の取り扱いと、その連続性。
- **C. 位相群**：反転 $g\mapsto g^{-1}$ の連続性と、合成（前/後）による連続性の保存。

\section*{記法と設定}
Lean 内の外部環境は概ね次のとおりです（型変数と位相・群の構造）:
\begin{itemize}
  \item 型 $T,X,Y,Z,W,G,H,K$、位相構造 $[{TopologicalSpace}\,\_\ ]$、群構造 $[{Group}\,\_\ ]$。
  \item 位相群は \verb|IsTopologicalGroup| として与えられ、積・反転の連続性を仮定します。
  \item 射影は $\pi_1:X\times Y\to X,\ \pi_2:X\times Y\to Y$。
\end{itemize}

\section*{A. 積位相}
### A. 積位相と射影

積位相の**普遍性**によれば、$X\times Y$ への関数 $h:T\to X\times Y$ が連続であることは、
二つの射影 $\pi_1:X\times Y\to X,\ \pi_2:X\times Y\to Y$ によって
$$
\pi_1\circ h : T\to X,\qquad \pi_2\circ h : T\to Y
$$
がいずれも連続であることと**同値**です。
Lean 側では、次のような定理がそれを表現しています：

- `continuous_π₁, continuous_π₂`：$\pi_1,\pi_2$ 自身が連続。
- `continuous_iff_proj_core, continuous_iff_proj, continuous_of_proj`：
  $h$ の連続性 $\Leftrightarrow$ 射影合成の連続性（必要十分・片方向）が明示されます。
- `continuous_pair_mk, continuous_Prod_map`：
  連続写像 $f:X\to Z,\ g:Y\to W$ から積写像
  $$
  f\times g : X\times Y \longrightarrow Z\times W,\quad (x,y)\mapsto (f(x),g(y))
  $$
  が連続であることを主張します。
- `prod_map_comp, prod_map_id_id`：`Prod.map` の合成公式と単位律
  $$
  \mathrm{Prod.map}(f_2,g_2)\circ \mathrm{Prod.map}(f_1,g_1)
   \;=\; \mathrm{Prod.map}(f_2\circ f_1,\ g_2\circ g_1),
  \qquad
  \mathrm{Prod.map}(\mathrm{id},\mathrm{id})=\mathrm{id}
  $$
  を確認します。

**証明の枠組み（概略）**：積位相は「射影が連続になるような最弱の位相」で、
この初期（イニシャル）構造の特徴づけから上の同値が従います。
$f\times g$ の連続性は、$\pi_1\circ(f\times g)=f\circ\pi_1$、
$\pi_2\circ(f\times g)=g\circ\pi_2$ の可換性により射影の連続性へ帰着します。

\subsection*{図式：射影と積写像}

\begin{tikzcd}
T \arrow[r, "h"] \arrow[dr, swap, "\pi_1\circ h"] \arrow[drr, "\pi_2\circ h"] & X\times Y \arrow[d, swap, "\pi_1"] \arrow[dr, "\pi_2"] \\
& X & Y
\end{tikzcd}


\medskip


\begin{tikzcd}
X\times Y \arrow[r, "{\mathrm{Prod.map}(f,g)}"] \arrow[d, swap, "\pi_1"] \arrow[dr, dashed, "{(f\circ\pi_1,\ g\circ\pi_2)}"] & Z\times W \arrow[d, "\pi_1"] \\
X \arrow[r, swap, "f"] & Z
\end{tikzcd}
\qquad
\begin{tikzcd}
X\times Y \arrow[r, "{\mathrm{Prod.map}(f,g)}"] \arrow[d, swap, "\pi_2"] \arrow[dr, dashed, "{(f\circ\pi_1,\ g\circ\pi_2)}"] & Z\times W \arrow[d, "\pi_2"] \\
Y \arrow[r, swap, "g"] & W
\end{tikzcd}


\section*{B. 群準同型の合成}
### B. 群準同型の合成

群（実際には `G →* H` という **MonoidHom**）の合成
$$
G \xrightarrow{\ \varphi\ } H \xrightarrow{\ \psi\ } K
\quad\leadsto\quad
G \xrightarrow{\ \psi\circ\varphi\ } K
$$
を Lean では `groupHomComp` として定義しています。

- `groupHomComp`：$\psi\circ\varphi$ を群準同型として与えます。
- `continuous_groupHomComp`：$\varphi,\psi$ が連続なら、その合成も連続。

**直観**：合成の連続性は一般の位相空間の写像でも成り立つ基本法則
（連続写像の合成は連続）にほぼ尽きます。

\subsection*{図式：合成の可視化}

\begin{tikzcd}
G \arrow[r, "\varphi"] \arrow[rr, bend left=20, "{\psi\circ\varphi}"] & H \arrow[r, "\psi"] & K
\end{tikzcd}


\section*{C. 位相群と反転の連続性}
### C. 位相群と反転の連続性

位相群 $(G,\tau_G)$ とは、群演算 $g\mapsto g^{-1}$ と $(g,h)\mapsto gh$
が**連続**となるような位相をもつ群です（Lean では `IsTopologicalGroup`）。
この設定下で、群準同型 $f:G\to H$ が連続なら、次も連続です：

- `topological_group_hom_continuous_precomp_inv`：$g\mapsto f(g^{-1})$（**前合成**での反転）。
- `topological_group_hom_continuous_postcomp_inv`：$g\mapsto (f(g))^{-1}$（**後合成**での反転）。
- `topological_group_hom_continuous_inv`：上と同値な主張の総括。

**証明の枠組み（概略）**：
$g\mapsto g^{-1}$ は位相群の公理により連続、$f$ も連続、
したがって合成も連続です。後合成の方は、$H$ 側の反転が連続であることを使います。

\subsection*{図式：反転と合成の連続性}

\begin{tikzcd}
G \arrow[r, "f"] \arrow[d, swap, "{\mathrm{inv}}"] & H \arrow[d, "{\mathrm{inv}}"] \\
G \arrow[r, swap, "f"] & H
\end{tikzcd}
\quad
\text{（縦の反転が連続 $\Rightarrow$ 矢印の合成も連続）}


\section*{Lean 内の定理・定義（一覧）}
\begin{center}
\begin{tabular}{lll}
\textbf{種別} & \textbf{名前} & \textbf{シグネチャ（先頭行）} \\\hline
theorem & \texttt{continuous\_iff\_proj\_core} & \texttt{{h : T → X × Y} : Continuous h ↔ (Continuous fun t => (h t).1) ∧ (Continuous fun t => (h t).2) := by} \\
def & \texttt{π₁} & \texttt{: X × Y → X := fun p => p.1} \\
def & \texttt{π₂} & \texttt{: X × Y → Y := fun p => p.2} \\
lemma & \texttt{continuous\_π₁} & \texttt{: Continuous (π₁ : X × Y → X) := by simpa using continuous\_fst} \\
theorem & \texttt{continuous\_pair\_mk} & \texttt{{f : X → Z} {g : Y → W} (hf : Continuous f) (hg : Continuous g) : Continuous (fun p : X × Y => (f p.1, g p.2)) := by} \\
theorem & \texttt{continuous\_iff\_proj} & \texttt{{h : T → X × Y} : Continuous h ↔ (Continuous fun t => (h t).1) ∧ (Continuous fun t => (h t).2) := by} \\
theorem & \texttt{continuous\_of\_proj} & \texttt{{h : T → X × Y} (h₁ : Continuous fun t => (h t).1) (h₂ : Continuous fun t => (h t).2) : Continuous h :=} \\
theorem & \texttt{continuous\_proj\_fst} & \texttt{{h : T → X × Y} (hh : Continuous h) : Continuous (fun t => (h t).1) :=} \\
theorem & \texttt{continuous\_proj\_snd} & \texttt{{h : T → X × Y} (hh : Continuous h) : Continuous (fun t => (h t).2) :=} \\
theorem & \texttt{continuous\_Prod\_map} & \texttt{{X Y Z W : Type*} [TopologicalSpace X] [TopologicalSpace Y] [TopologicalSpace Z] [TopologicalSpace W] {f : X → Z} {g : Y → W} (hf : Continuous f) (hg : Continuous g) : Continuous (Prod.map f g) := by} \\
theorem & \texttt{prod\_map\_comp} & \texttt{ {A B C D E F : Type*} (f₁ : A → B) (g₁ : C → D) (f₂ : B → E) (g₂ : D → F) : (Prod.map f₂ g₂) ∘ (Prod.map f₁ g₁) = (Prod.map (f₂ ∘ f₁) (g₂ ∘ g₁)) := by} \\
theorem & \texttt{prod\_map\_id\_id} & \texttt{{X Y : Type*} : Prod.map (id : X → X) (id : Y → Y) = (id : X × Y → X × Y) := by} \\
def & \texttt{groupHomComp} & \texttt{(φ : G →* H) (ψ : H →* K) : G →* K := ψ.comp φ} \\
theorem & \texttt{continuous\_groupHomComp} & \texttt{ (φ : G →* H) (ψ : H →* K) (hf : Continuous φ) (hg : Continuous ψ) : Continuous (fun g : G => groupHomComp φ ψ g) := by} \\
theorem & \texttt{topological\_group\_hom\_continuous\_precomp\_inv} & \texttt{ (f : G →* H) (hf : Continuous f) : Continuous (fun g : G => f (g⁻¹)) := by} \\
theorem & \texttt{topological\_group\_hom\_continuous\_postcomp\_inv} & \texttt{ (f : G →* H) (hf : Continuous f) : Continuous (fun g : G => (f g)⁻¹) := by} \\
theorem & \texttt{topological\_group\_hom\_continuous\_inv} & \texttt{ (f : G →* H) (hf : Continuous f) : Continuous (fun g : G => (f g)⁻¹) :=} \\
\end{tabular}
\end{center}

\section*{Lean ソース（全文）}
\begin{lstlisting}[language=lean]
/-
  Bourbaki-style: ordered pairs, projections, and morphisms

  This module follows the spirit of Nicolas Bourbaki's Elements of Mathematics:
  structures carried by sets, ordered pairs characterized via projections (π₁, π₂),
  and morphisms checked by their behavior under these projections.

  Contents (Tasks A–C):
  A. Product topology — continuity of the induced product map (f × g)
  B. Group homomorphisms — composition as a morphism in the category of groups
  C. Topological groups — continuity under inversion (pre/post composition)
-/

import Mathlib.Topology.Basic
import Mathlib.Topology.Constructions.SumProd
import Mathlib.Topology.Algebra.Group.Basic
import Mathlib.Algebra.Group.Basic
import Mathlib.Tactic
import Mathlib.Topology.ContinuousMap.Basic

noncomputable section

namespace MyProjects.Topology.C

open scoped Topology

-- Core equivalence placed early for use in subsequent lemmas.
section ProductUPCore

variable {T X Y : Type*}
variable [TopologicalSpace T] [TopologicalSpace X] [TopologicalSpace Y]

theorem continuous_iff_proj_core {h : T → X × Y} :
    Continuous h ↔ (Continuous fun t => (h t).1) ∧ (Continuous fun t => (h t).2) := by
  constructor
  · intro hh
    exact ⟨continuous_fst.comp hh, continuous_snd.comp hh⟩
  · intro hproj
    rcases hproj with ⟨h₁, h₂⟩
    simpa using h₁.prodMk h₂

end ProductUPCore

/-! # Ordered pairs and projections

In Bourbaki's viewpoint, the ordered pair (x, y) is determined by its projections
π₁, π₂. The product topology on X × Y is the coarsest topology making these
projections continuous, so continuity of maps into products can be checked
componentwise via π₁, π₂.
-/

section ProductTopology

variable {X Y Z W : Type*}
variable [TopologicalSpace X] [TopologicalSpace Y]
variable [TopologicalSpace Z] [TopologicalSpace W]

/-- First projection π₁ : X × Y → X. -/
def π₁ : X × Y → X := fun p => p.1

/-- Second projection π₂ : X × Y → Y. -/
def π₂ : X × Y → Y := fun p => p.2

@[simp] lemma π₁_apply (p : X × Y) : π₁ p = p.1 := rfl
@[simp] lemma π₂_apply (p : X × Y) : π₂ p = p.2 := rfl

lemma continuous_π₁ : Continuous (π₁ : X × Y → X) := by simpa using continuous_fst
lemma continuous_π₂ : Continuous (π₂ : X × Y → Y) := by simpa using continuous_snd

/--
Task A. If `f : X → Z` and `g : Y → W` are continuous, then the induced product
map `X × Y → Z × W, p ↦ (f p.1, g p.2)` is continuous. This is the universal
property of products expressed via projections.
-/
@[continuity]
theorem continuous_pair_mk {f : X → Z} {g : Y → W}
    (hf : Continuous f) (hg : Continuous g) :
    Continuous (fun p : X × Y => (f p.1, g p.2)) := by
  -- Use the universal property: continuity is equivalent to continuity of projections.
  have hfproj : Continuous (fun p : X × Y => (f p.1, g p.2).1) := by
    simpa using (hf.comp continuous_fst)
  have hgproj : Continuous (fun p : X × Y => (f p.1, g p.2).2) := by
    simpa using (hg.comp continuous_snd)
  exact (continuous_iff_proj_core (h := fun p : X × Y => (f p.1, g p.2))).2 ⟨hfproj, hgproj⟩

-- Quick sanity checks (do not export public API)
section Examples
  variable {f : X → Z} {g : Y → W}
  example (hf : Continuous f) (hg : Continuous g) :
      Continuous (fun p : X × Y => (f p.fst, g p.snd)) :=
    by simpa using continuous_pair_mk (X:=X) (Y:=Y) (Z:=Z) (W:=W) hf hg
  example : Continuous (fun p : X × Y => (p.1, p.2)) := by
    continuity
end Examples

end ProductTopology

/-! # Universal property via projections

For a map into a product, continuity is equivalent to continuity of its
compositions with the projections. This encodes Bourbaki’s “ordered pair is
read by its projections” principle as an equivalence.
-/

section ProductUP

variable {T X Y : Type*}
variable [TopologicalSpace T] [TopologicalSpace X] [TopologicalSpace Y]

/-- Continuity into a product is equivalent to continuity of both projections. -/
theorem continuous_iff_proj {h : T → X × Y} :
    Continuous h ↔ (Continuous fun t => (h t).1) ∧ (Continuous fun t => (h t).2) := by
  simpa using (continuous_iff_proj_core (h:=h))

/-- One-sided version: if both projections are continuous, then `h` is continuous.
    Kept separate to avoid `simp` loops with the equivalence. -/
theorem continuous_of_proj {h : T → X × Y}
    (h₁ : Continuous fun t => (h t).1) (h₂ : Continuous fun t => (h t).2) :
    Continuous h :=
  (continuous_iff_proj (h:=h)).2 ⟨h₁, h₂⟩

/-- One-sided version: continuity implies continuity of the first projection. -/
theorem continuous_proj_fst {h : T → X × Y}
    (hh : Continuous h) : Continuous (fun t => (h t).1) :=
  continuous_fst.comp hh

/-- One-sided version: continuity implies continuity of the second projection. -/
theorem continuous_proj_snd {h : T → X × Y}
    (hh : Continuous h) : Continuous (fun t => (h t).2) :=
  continuous_snd.comp hh

/-- Bridge for `Prod.map`: continuity from componentwise continuity. -/
@[continuity]
theorem continuous_Prod_map {X Y Z W : Type*}
    [TopologicalSpace X] [TopologicalSpace Y]
    [TopologicalSpace Z] [TopologicalSpace W]
    {f : X → Z} {g : Y → W}
    (hf : Continuous f) (hg : Continuous g) :
    Continuous (Prod.map f g) := by
  -- Use the universal property directly on `h := Prod.map f g`.
  apply (continuous_iff_proj_core (h := Prod.map f g)).2
  constructor
  · -- first projection continuous
    simpa [Prod.map] using (hf.comp continuous_fst)
  · -- second projection continuous
    simpa [Prod.map] using (hg.comp continuous_snd)

/-- Normalization: composition of `Prod.map` equals `Prod.map` of compositions. -/
@[simp]
theorem prod_map_comp
    {A B C D E F : Type*}
    (f₁ : A → B) (g₁ : C → D) (f₂ : B → E) (g₂ : D → F) :
    (Prod.map f₂ g₂) ∘ (Prod.map f₁ g₁) =
      (Prod.map (f₂ ∘ f₁) (g₂ ∘ g₁)) := by
  funext p; cases p with
  | _ a c =>
    simp [Function.comp, Prod.map]

/-- Normalization: `Prod.map` with identities is the identity. -/
@[simp]
theorem prod_map_id_id {X Y : Type*} :
    Prod.map (id : X → X) (id : Y → Y) = (id : X × Y → X × Y) := by
  funext p; cases p with
  | _ x y => simp [Prod.map]

end ProductUP

/-! # Triangle identities (componentwise)

Right-triangle style identities for products, stated componentwise as function
equalities. These are convenient `simp` targets and make `by continuity` flow
through compositions with projections and `Prod.map`.
-/

section ProductTriangles

variable {T X Y Z W : Type*}

/-- Componentwise identity: composing with the first projection undoes pairing. -/
@[simp] theorem fst_comp_pair (f : T → X) (g : T → Y) :
    (fun t => ((f t, g t)).1) = f := by
  funext t; rfl

/-- Componentwise identity: composing with the second projection undoes pairing. -/
@[simp] theorem snd_comp_pair (f : T → X) (g : T → Y) :
    (fun t => ((f t, g t)).2) = g := by
  funext t; rfl

/-- Componentwise normalization for `Prod.map` and `fst`. -/
@[simp] theorem fst_comp_Prod_map (f : X → Z) (g : Y → W) :
    (Prod.fst ∘ Prod.map f g) = (f ∘ Prod.fst) := by
  funext p; cases p with
  | _ x y => simp [Function.comp, Prod.map]

/-- Componentwise normalization for `Prod.map` and `snd`. -/
@[simp] theorem snd_comp_Prod_map (f : X → Z) (g : Y → W) :
    (Prod.snd ∘ Prod.map f g) = (g ∘ Prod.snd) := by
  funext p; cases p with
  | _ x y => simp [Function.comp, Prod.map]

end ProductTriangles

/-! # Group morphisms

In Bourbaki's structural approach, morphisms between groups compose to yield a
group morphism again; categorically, this is just composition of arrows.
-/

section GroupHom

variable {G H K : Type*}
variable [Group G] [Group H] [Group K]

/-- Task B. Composition of group homomorphisms (implemented as `MonoidHom`). -/
def groupHomComp (φ : G →* H) (ψ : H →* K) : G →* K :=
  ψ.comp φ

@[simp] lemma groupHomComp_apply (φ : G →* H) (ψ : H →* K) (g : G) :
    groupHomComp φ ψ g = ψ (φ g) := rfl

@[simp] lemma groupHomComp_eq_comp (φ : G →* H) (ψ : H →* K) :
    groupHomComp φ ψ = ψ.comp φ := rfl

-- Left identity (id on the domain of ψ): ψ ∘ id_G = ψ
@[simp] lemma groupHomComp_id_left (ψ : G →* K) :
    groupHomComp (G:=G) (H:=G) (K:=K) (MonoidHom.id G) ψ = ψ := by
  ext g; simp [groupHomComp]

-- Variant with `G := H` (often triggered by type inference on `ψ : H →* K`).
@[simp] lemma groupHomComp_id_left' (ψ : H →* K) :
    groupHomComp (G:=H) (H:=H) (K:=K) (MonoidHom.id H) ψ = ψ := by
  ext g; simp [groupHomComp]

-- Right identity specialized (codomain equals H): ψ = id_H ∘ φ
@[simp] lemma groupHomComp_id_right (φ : G →* H) :
    groupHomComp (G:=G) (H:=H) (K:=H) φ (MonoidHom.id H) = φ := by
  ext g; simp [groupHomComp]

section TopologyBridge
  variable [TopologicalSpace G] [TopologicalSpace H] [TopologicalSpace K]
  /-- Bridge: continuity of the underlying function of `groupHomComp`. -/
  @[continuity]
  theorem continuous_groupHomComp
      (φ : G →* H) (ψ : H →* K)
      (hf : Continuous φ) (hg : Continuous ψ) :
      Continuous (fun g : G => groupHomComp φ ψ g) := by
    simpa [groupHomComp] using (hg.comp hf : Continuous fun g : G => ψ (φ g))

  -- doctest: continuity tactic uses the bridge
  example (φ : G →* H) (ψ : H →* K)
      (hf : Continuous φ) (hg : Continuous ψ) :
      Continuous fun g : G => ψ (φ g) := by
    continuity
end TopologyBridge

-- Example: pointwise simplification
section Examples
  variable (φ : G →* H) (ψ : H →* K) (g : G)
  example : groupHomComp φ ψ g = ψ (φ g) := by simpa using groupHomComp_apply φ ψ g
end Examples

end GroupHom

/-! # Topological groups and inversion

Two Bourbaki-flavoured continuity principles around inversion for a continuous
homomorphism `f : G →* H` between topological groups.
-/

section TopologicalGroup

variable {G H : Type*}
variable [TopologicalSpace G] [Group G]
variable [TopologicalSpace H] [Group H]

section PrecomposeInv
  variable [IsTopologicalGroup G]
  /-- Pre-compose with inversion on the domain. -/
  @[continuity]
  theorem topological_group_hom_continuous_precomp_inv
      (f : G →* H) (hf : Continuous f) :
      Continuous (fun g : G => f (g⁻¹)) := by
    exact hf.comp (continuous_inv : Continuous fun g : G => g⁻¹)
end PrecomposeInv

section PostcomposeInv
  variable [IsTopologicalGroup H]
  /-- Post-compose with inversion on the codomain. -/
  @[continuity]
  theorem topological_group_hom_continuous_postcomp_inv
      (f : G →* H) (hf : Continuous f) :
      Continuous (fun g : G => (f g)⁻¹) := by
    exact (continuous_inv.comp hf : Continuous fun g : G => (f g)⁻¹)

  /-- Compatibility alias with earlier naming in notes. -/
  theorem topological_group_hom_continuous_inv
      (f : G →* H) (hf : Continuous f) :
      Continuous (fun g : G => (f g)⁻¹) :=
    topological_group_hom_continuous_postcomp_inv f hf

  -- Minimal examples to exercise the API (remain local to this section)
  section Examples
    variable (f : G →* H)
    example (hf : Continuous f) : Continuous (fun g : G => (f g)⁻¹) :=
      topological_group_hom_continuous_postcomp_inv f hf
    -- continuity tactic doctest
    example (hf : Continuous f) : Continuous fun g : G => (f g)⁻¹ := by
      continuity
  end Examples
end PostcomposeInv

end TopologicalGroup

end MyProjects.Topology.C

\end{lstlisting}

\end{document}
